\chapter*{Programação}
\addcontentsline{toc}{chapter}{Programação}


\section{Cronograma}
\addcontentsline{toc}{section}{Cronograma}

O cronograma abaixo detalha a distribuição das atividades ao longo dos 12 meses previstos para a execução do projeto.

\begin{table}[h!]
\centering
\caption{Cronograma de Execução das Atividades}
\label{tab:cronograma}
\resizebox{\textwidth}{!}{%
\begin{tabular}{|l|c|c|c|c|c|c|c|c|c|c|c|c|}
\hline
\multicolumn{1}{|c|}{\textbf{Atividades}} & \textbf{Fev} & \textbf{Mar} & \textbf{Abr} & \textbf{Mai} & \textbf{Jun} & \textbf{Ago} & \textbf{Set} & \textbf{Out} & \textbf{Nov} \\ \hline

1. Revisão Bibliográfica Sistemática e Fundamentação Teórica 
& x & x & & & & & & & \\ \hline

2. Fase 1 – Arquitetura de Visão e Aquisição de Dados 
& & x & x & & & & & & \\ \hline

3. Fase 2 – Modelagem, Impressão 3D e Montagem da Base Cinemática (2 DoF) 
& & & x & x & & & & & \\ \hline

4. Fase 3 – Coleta e Curadoria de Dataset 
& & & & x & x & & & & \\ \hline

5. Fase 4 – Treinamento, Otimização e Inferência Embarcada (Edge AI) 
& & & & x & x & x & & & \\ \hline

6. Fase 5 – Integração Visão–Controle e Cinemática Inversa 
& & & & & x & x & x & & \\ \hline

7. Fase 6 – Testes de Validação e Coleta de Métricas Experimentais 
& & & & & & x & x & & \\ \hline

8. Análise dos Resultados e Escrita de Artigos Científicos 
& & & & & & & x & x & x \\ \hline

9. Redação Final do Relatório Técnico / Projeto de IC 
& & & & & & & & x & x \\ \hline

\end{tabular}%
}
\end{table}

\section{Encontros de Acompanhamento}
\addcontentsline{toc}{section}{Encontros de Acompanhamento}

Os encontros de acompanhamento entre os orientadores e os alunos serão realizados semanalmente. Esses encontros terão como objetivo discutir o progresso das atividades, resolver dúvidas, ajustar o cronograma conforme necessário e garantir que os objetivos do projeto estejam sendo alcançados de forma eficaz.

Previsão de encontros:
\textbf{Segundas-feiras}, das \textbf{13:20h às 14:20h}.